\documentclass[../../main.tex]{subfiles}
\graphicspath{{../../image/}{../../figure/}}
\begin{document}

\section{强波的稳定性}\label{section:ch14-强波的稳定性}

本节证明标准 Riemann 半群轨道中强波的稳定性。其思想是把强激波（而非弱激波）视为解骨架的永久构件，在其上叠加小扰动。

\begin{theorem}\label{theorem:ch14-强波稳定性}
设 $U$ 是 $S_t$ 的轨道，初值 $U_0\in\mathscr{D}$，并设 $x=\xi(t)$ 是 $U$ 中任一强度 $|\gamma_i|>\bar\gamma$ 的 $i$-激波轨迹，其中 $\bar\gamma>0$ 充分小。则存在常数 $\omega>0$ 与 $C>0$，使对任一其他轨道 $\bar U$，其初值 $\bar U_0$ 与 $U_0$ 在强激波 $x=\xi(0)$ 附近一致，
\begin{align}
\|\bar U(\cdot,t)-U(\cdot,t)\|_{L^1}\leqslant C\|U_0-\bar U_0\|_{L^1}e^{\omega t},\qquad t\geqslant0.
\label{align:ch14-strong-wave-stability}
\end{align}
此外，$\bar U$ 中相应 $i$-激波的位置 $x=\bar\xi(t)$ 满足
\begin{align}
|\bar\xi(t)-\xi(t)|\leqslant C\|U_0-\bar U_0\|_{L^1},\qquad t\geqslant0.
\label{align:ch14-strong-wave-position}
\end{align}
\end{theorem}

\begin{proof}
我们通过适配\cref{section:ch14-连续依赖性}中建立的 $\rho$ 泛函来证明。固定强度 $|\gamma_i|>\bar\gamma$ 的强 $i$-激波 $x=\xi(t)$。此激波把上半平面分成左右两半，我们在每一半上独立测量 $U$ 与 $\bar U$ 间的距离。

在激波左侧，$U$ 与 $\bar U$ 的状态以距离 $\rho_L$ 测量；右侧以距离 $\rho_R$ 测量。强激波本身作为界面，其位置失配以 $|\bar\xi(t)-\xi(t)|$ 度量。

由\cref{section:ch14-连续依赖性}的论证，在激波左侧与右侧，$\rho_L$ 与 $\rho_R$ 满足形如\eqref{align:ch14-rho-decay}的微分不等式。此外，强激波的位置失配 $|\bar\xi(t)-\xi(t)|$ 由 Rankine--Hugoniot 条件控制：由于激波强度 $|\gamma_i|>\bar\gamma$，激波速度对状态扰动的敏感度有正常数下界，因而位置失配在时间演化中不指数增长。

具体地，令 $E(t)=\rho_L(t)+\rho_R(t)+|\bar\xi(t)-\xi(t)|$。由波相互作用分析可得
\begin{align}
\frac{\mathrm{d}}{\mathrm{d}t}E(t)\leqslant\omega E(t),
\label{align:ch14-E-diff}
\end{align}
其中 $\omega$ 仅依赖于强激波强度与 $F$ 的导数界。Gronwall 不等式给出\eqref{align:ch14-strong-wave-stability}。

\eqref{align:ch14-strong-wave-position}由激波速度对状态失配的 Lipschitz 连续得到：由于激波速度差以状态 $L^1$ 距离为界，位置失配在时间上累积但不超过初始 $L^1$ 距离乘以常数 $C$。
\end{proof}

\cref{theorem:ch14-强波稳定性}表明，强激波一旦形成便在小扰动下稳定，其位置仅以初始扰动 $L^1$ 距离为界偏移。这与弱激波不同，后者在波相互作用中可能被吸收或分裂。

强波稳定性的另一重要推论是解的``结构刚性''：在标准 Riemann 半群轨道中，强激波的数目与位置在小扰动下保持不变，而弱波在激波之间传播并最终相互作用。

\end{document}
